% Figure: necklaces for the FLT counting proof --- p beads, a colours.
% Source for Figure~\ref{fig:necklace}. TikZ only; regenerable by compiling main.tex.
\begin{tikzpicture}[
  bead/.style={circle,draw=ink,minimum size=5mm,font=\scriptsize},
  mono/.style={bead,fill=accentsoft},
  lab/.style={font=\footnotesize\sffamily,text=slate},
]
% 5-bead necklace, pattern A B A A C
\foreach \i/\c in {0/A,1/B,2/A,3/A,4/C}
  \node[bead] at ({90+\i*72}:1.1) {\c};
\node[lab] at (0,-1.9) {one necklace, $p = 5$ beads, $a$ colour names};
% three monochrome necklaces (small rings at x = 3.0, 4.3, 5.6)
\foreach \k in {0,1,2}
{
  \foreach \i in {0,...,4}
  {
    \pgfmathsetmacro{\px}{3.0 + \k*1.3 + 0.42*cos(90+\i*72)}
    \pgfmathsetmacro{\py}{0.42*sin(90+\i*72)}
    \node[mono] at (\px,\py) {};
  }
  \node[lab] at (3.0+\k*1.3,-1.0) {\ifnum\k=0 all-A\else\ifnum\k=1 all-B\else all-C\fi\fi};
}
\node[lab] at (4.3,-1.9) {$a$ monochrome necklaces};
\node[lab] at (0,-2.7) {the other $a^p - a$ strings sit in rotation-classes of size $p$};
\node[lab] at (0,-3.1) {($p$ prime: an orbit smaller than $p$ would repeat a proper divisor)};
\node[lab] at (0,-3.5) {$\Rightarrow\; a^p - a = p\cdot(\text{classes}) \equiv 0 \pmod p$};
\end{tikzpicture}
